\chapter{Piret--Quisquater Differential Fault Analysis of AES}
\section{Where the fault is injected}
In the classic model used in the notes, one nonzero byte is corrupted at the \emph{input of MixColumns in round 9}, after SubBytes and ShiftRows. Both correct and faulty executions then proceed through round 9 and the final round.
\begin{center}\begin{tikzpicture}[node distance=5mm, every node/.style={font=\small}, box/.style={draw,rounded corners,fill=pale,minimum width=8cm,minimum height=8mm},>=Latex]
\node[box] (a) {Round 9: SubBytes $\to$ ShiftRows};
\node[box,below=of a] (b) {Inject one-byte XOR fault};
\node[box,below=of b] (c) {Round 9: MixColumns $\to$ AddRoundKey};
\node[box,below=of c] (d) {Round 10: SubBytes $\to$ ShiftRows $\to$ AddRoundKey};
\node[box,below=of d] (e) {Compare correct $C$ and faulty $F$};
\foreach \x/\y in {a/b,b/c,c/d,d/e}{\draw[->] (\x)--(\y);}
\end{tikzpicture}\end{center}
\section{Propagation through MixColumns}
If the first byte of a column changes by \(e\), the column difference after MixColumns is
\[\Delta_{MC}=(02e,\ e,\ e,\ 03e)^T.\]
For \(e=\hex{01}\), this is \((\hex{02},\hex{01},\hex{01},\hex{03})^T\). AddRoundKey does not alter XOR differences because both executions use the same round key. The final SubBytes transforms these differences nonlinearly; ShiftRows moves them into separate ciphertext positions.
\section{Which ciphertext bytes are affected?}
For the four round-9 columns, the affected ciphertext positions (in corresponding row order) are:
\begin{center}\begin{tabular}{cc}\toprule Faulted column & Ciphertext byte positions\\\midrule 0&0, 13, 10, 7\\1&4, 1, 14, 11\\2&8, 5, 2, 15\\3&12, 9, 6, 3\\\bottomrule\end{tabular}\end{center}
\section{Constructing candidates for the last AES round key}
For an affected ciphertext position \(j\), and a guessed last-round key byte \(k\), compute
\[d_j(k)=S^{-1}(C_j\xor k)\xor S^{-1}(F_j\xor k).\]
Each byte has 256 candidate values. Four unknown key bytes have \(256^4=2^{32}\) possible combinations before filtering.
Construct a set \(\mathcal D\) containing the MixColumns output differences for a single nonzero faulty byte in any of the four row positions:
\[\mathcal D=\{M(e\mathbf e_r):r\in\{0,1,2,3\},\ e\in\{1,\ldots,255\}\}.\]
This produces \(4\cdot255=1020\) four-byte difference patterns (not 1020 individual bytes). For each proposed key quartet, compute the four \(d_j(k)\) values in their correct \emph{pre-ShiftRows row order} and keep the quartet only if the resulting vector lies in \(\mathcal D\).
\section{Intersecting observations}
A single correct/faulty pair typically leaves multiple candidate key quartets. Repeat with independent faults in the same column and intersect the candidate sets:
\[\mathcal K_{\mathrm{surviving}}=\mathcal K_1\cap\mathcal K_2\cap\cdots.\]
Repeat across all four columns to recover the full 16-byte \(K_{10}\). The notes illustrate candidate filtering using a hypothetical quartet \((\hex{13},\hex{2B},\hex{A7},\hex{17})\); such pedagogical candidate examples should not be confused with a complete independently verified AES fault trace.
\section{Two injection depths: round 9 versus round 8}
\label{sec:piret_two_models}
The same Piret principle can be applied at two different depths, and the project's
analysis distinguishes them because they behave very differently. The choice is
the \texttt{fault\_round} argument of the target in \texttt{piret\_aes.py}.

\begin{description}
\item[Round-9 model (one column group).] The fault is injected before the
round-9 MixColumns, as described above. MixColumns spreads the single faulted
byte over \emph{one} column (four bytes); ShiftRows then scatters those four
bytes to four ciphertext positions. One faulty ciphertext therefore constrains
only \emph{one} group of four last-round key bytes. To recover all sixteen bytes
of \(K_{10}\), the campaign must be repeated for each of the four columns.
\item[Round-8 model (all four groups at once).] The fault is injected one round
earlier, before the round-8 MixColumns. Now the difference passes through
\emph{two} MixColumns layers: round-8 MixColumns spreads the single byte to a
full column, and round-9 MixColumns spreads each of those bytes across all four
columns. Before the final round the difference already touches all sixteen state
bytes, so a single faulty ciphertext constrains \emph{all four} key-byte groups
simultaneously. This is why the measured campaigns recover the complete
last-round key from a median of only two faulty ciphertexts.
\end{description}

\noindent The trade-off is depth versus count: the round-8 model needs far fewer
faulty ciphertexts but assumes a cleaner, more deeply propagated fault and larger
per-observation candidate structures; the round-9 model is simpler to reason
about per column but must be repeated. The default command line and the analysis
use the round-8 model; \texttt{run.py} follows the round-9 model used in the
lecture notes. Both reduce to the same last-round key once enough observations
are intersected.

\section{Reversing AES-128 key expansion}
Let \(w_0,\ldots,w_{43}\) be 32-bit expanded-key words. For \(i\ge4\),
\[w_i=\begin{cases}w_{i-4}\xor g(w_{i-1}),&i\equiv0\pmod4,\\w_{i-4}\xor w_{i-1},&\text{otherwise},\end{cases}\]
where \(g(w)=\operatorname{SubWord}(\operatorname{RotWord}(w))\xor\operatorname{Rcon}\). In reverse, compute \(w_{i-4}\) by XORing \(w_i\) with the same known right-hand transformation. Begin with \(w_{40},\ldots,w_{43}\) from \(K_{10}\) and work backward to \(w_0,\ldots,w_3\), yielding \(K_0\). No inverse S-box is required for this reversal because the relevant inputs to \(g\) are already known.
